%=============================================================================!
% 14.4  THE PRESSURE TENSOR                                                   !
%=============================================================================!
\section{The pressure tensor}
\label{sec:pr-tensor}

A liquid at rest pushes equally on a face whichever way the face points.
A solid need not: a crystal stretched along $x$ pulls inwards harder on
the faces normal to $x$ than on the others, and a sheet of graphite is
stiff along its plane and soft across it. This section generalises the
pressure to a matrix that gives the push on a face of any orientation.

\subsection*{The push on a face}

The force that the material exerts on a small flat face of area $A$,
whose outward normal is the unit vector $\vb{n}$, need not point along
$\vb{n}$, but it is linear in $\vb{n}$. A small wedge of material cut off
by this face and by three faces normal to the axes, of areas $An_x$,
$An_y$ and $An_z$, feels forces on its faces that shrink with their area,
while its mass shrinks faster, with its volume; as the wedge shrinks the
forces on its four faces must therefore cancel, and the force on the
slanted face is the sum of $An_x$, $An_y$ and $An_z$ times the forces per
unit area on faces normal to $x$, $y$ and $z$. A vector linear in $\vb{n}$
is a matrix times $\vb{n}$ (\cref{sec:os-modes}):
\begin{equation*}
   \vb{F} = A\,\mathsf{P}\vb{n} ,
\end{equation*}
with $\mathsf{P}$ the pressure tensor, whose component
$\mathsf{P}_{\alpha\beta}$ is the force per unit area along $\alpha$ on a
face normal to $\beta$. For a liquid at rest $\mathsf{P} = P\vb{I}$, and
every face is pushed straight outwards with $P$.

\subsection*{From a strain}

\Cref{sec:pr-periodic} stretched the box evenly; stretching it unevenly
gives every component. A small strain multiplies every vector, the
lattice vectors of the cell and the position of every atom alike, by the
matrix $\vb{I} + \boldsymbol{\epsilon}$, with components
$\epsilon_{\alpha\beta}$ that are small. For a cubic cell of side $L$ the
strained lattice vectors are $L$ times the columns $\vb{c}_1$, $\vb{c}_2$
and $\vb{c}_3$ of $\vb{I} + \boldsymbol{\epsilon}$, so by
\cref{eq:md-volume} the volume becomes $V$ times the triple product
$\vb{c}_1\cdot(\vb{c}_2\times\vb{c}_3)$, which is the determinant
$\det(\vb{I} + \boldsymbol{\epsilon})$, the factor by which the strain
stretches volumes (\cref{sec:ha-liouville}). Written out as in
\cref{sec:ro-tensor}, it is six products, each of one component of every
column, and to first order in $\boldsymbol{\epsilon}$
\begin{equation*}
   \det(\vb{I} + \boldsymbol{\epsilon}) = (1 + \epsilon_{xx})(1 +
   \epsilon_{yy})(1 + \epsilon_{zz}) + \dotsb = 1 + \epsilon_{xx} +
   \epsilon_{yy} + \epsilon_{zz} + \dotsb ,
\end{equation*}
since each of the other five products contains at least two off-diagonal
components and is of second order or higher (\cref{ex:pr-det}).

The strain $\dd\epsilon_{\alpha\beta}$ carries the face of a cubic box
normal to $\beta$, at the distance $L$ from the opposite face, a distance
$L\,\dd\epsilon_{\alpha\beta}$ along $\alpha$ further than that opposite
face. The material pushes on the face with $L^2\mathsf{P}_{\alpha\beta}$
along $\alpha$, so it does the work $V\mathsf{P}_{\alpha\beta}\,
\dd\epsilon_{\alpha\beta}$, which its free energy loses at fixed
temperature:
\begin{equation*}
   \dd F = -V\sum_{\alpha\beta}\mathsf{P}_{\alpha\beta}\,
   \dd\epsilon_{\alpha\beta} ,
\end{equation*}
and so $\mathsf{P}_{\alpha\beta} = -(1/V)\,\partial F/\partial
\epsilon_{\alpha\beta}$; an even stretch, $\epsilon_{\alpha\beta} =
e\,\delta_{\alpha\beta}$, changes the volume by $\dd V = 3V\dd e$ and gives
back $\dd F = -P\,\dd V$ with $P$ the mean of the diagonal. The steps of
\cref{sec:pr-periodic} carry over. In fractional coordinates the
partition function carries the factor $\det(\vb{I} +
\boldsymbol{\epsilon})^N$, whose logarithm, $N\ln(1 + \epsilon_{xx} +
\epsilon_{yy} + \epsilon_{zz} + \dotsb)$, has the slope
$N\delta_{\alpha\beta}$ along $\epsilon_{\alpha\beta}$ at zero strain. Each
pair separation $\vb{d}$, of length $r_{ij}$, becomes $(\vb{I} +
\boldsymbol{\epsilon})\vb{d}$, with components $d_\alpha +
\sum_\beta\epsilon_{\alpha\beta}d_\beta$; squaring and adding them, and
dropping the products of two strains, gives the squared length $r_{ij}^2 +
2\sum_{\alpha\beta}d_\alpha\epsilon_{\alpha\beta}d_\beta$, and since
$\sqrt{r^2 + x} \approx r + x/2r$ for small $x$ (\cref{sec:mo-taylor}),
the length grows by $\sum_{\alpha\beta}d_\alpha\epsilon_{\alpha\beta}
d_\beta/r_{ij}$. Its slope along $\epsilon_{\alpha\beta}$ is $d_\alpha
d_\beta/r_{ij}$, and by the chain rule
\begin{equation*}
   \frac{\partial U}{\partial\epsilon_{\alpha\beta}}
   = \sum_{i<j}\varphi'(r_{ij})\,\frac{d_\alpha d_\beta}{r_{ij}}
   = -\mathcal{W}_{\alpha\beta} ,
\end{equation*}
with $\mathcal{W}_{\alpha\beta}$ the virial tensor, whose trace, the sum of
its diagonal, is $\mathcal{W}$. Then, as in \cref{sec:pr-periodic},
\begin{equation*}
   \mathsf{P}_{\alpha\beta} = \frac{\kB T}{V}\,
   \frac{\partial\ln Z}{\partial\epsilon_{\alpha\beta}}
   = \frac1V\bigl(N\kB T\,\delta_{\alpha\beta}
   + \ens{\mathcal{W}_{\alpha\beta}}\bigr) .
\end{equation*}
Equipartition gives $\ens{m_iv_{i\alpha}^2} = \kB T$, and in the Maxwell
distribution (\cref{sec:sm-maxwell}) the components of a velocity are
independent with mean zero, so $\ens{m_iv_{i\alpha}v_{i\beta}} = \kB T\,
\delta_{\alpha\beta}$, and the instantaneous
\begin{equation}
   \mathsf{P}_{\alpha\beta} = \frac1V\Bigl(\sum_im_iv_{i\alpha}v_{i\beta}
   + \mathcal{W}_{\alpha\beta}\Bigr)
   \label{eq:pr-tensor}
\end{equation}
averages to the pressure tensor, and the mean of its diagonal is the
instantaneous pressure of \cref{sec:pr-virial}. \texttt{md.PairModel}
stores $\mathcal{W}_{\alpha\beta}$ at every call, from the separations
$\vb{d}$ of the pairs and the forces between them:

\begin{code}
        r = np.asarray(positions, dtype=float)
        i, j, d, dist = self.neighbours.pairs(r)
        phi, dphi = self._pair(i, j, dist)
        push = (dphi / dist)[:, None] * d  # the force on i from j
        forces = np.zeros_like(r)
        np.add.at(forces, i, push)
        np.add.at(forces, j, -push)
        self.virial = -np.einsum("pa,pb->ab", d, push)
        return float(np.sum(phi)), forces
\end{code}

\noindent \texttt{d} holds $\vb{r}_j - \vb{r}_i$ to the nearest copy and
\texttt{push} the force on $i$ from $j$, $\varphi'(r_{ij})\,\vb{d}/r_{ij}$
(\cref{sec:pe-forces}), so \texttt{virial} is $-\sum d_\alpha F_\beta =
-\sum\varphi'(r_{ij})\,d_\alpha d_\beta/r_{ij}$ over the pairs,
$\mathcal{W}_{\alpha\beta}$. It agrees with $-\partial U/\partial
\epsilon_{\alpha\beta}$ for each of the nine strains, taken by finite
differences, to \SI{2.1e-9}{\electronvolt}, and $-\mathcal{W}_{\alpha\beta}
/V$ agrees with the stress of ASE's \texttt{LennardJones} to
\SI{7.5e-16}{\giga\pascal}.

\subsection*{Measured}

The liquid's tensor is isotropic: under the cell rescaling of
\cref{sec:pr-rescaling} at \SI{0.1}{\giga\pascal} its diagonal averages
0.1003, 0.0995 and \SI{0.0996}{\giga\pascal} and its off-diagonal
components $-0.0002$, 0.0000 and \SI{-0.0001}{\giga\pascal}, while an
instantaneous diagonal component scatters by \SI{0.023}{\giga\pascal} and
an off-diagonal one by \SI{0.013}{\giga\pascal}. The crystal at $a_0$
stretched along $x$ is not (\cref{fig:pr-crystal}(b)): at
$\epsilon_{xx} = +0.02$, $\mathsf{P}_{xx} = \SI{-0.0735}{\giga\pascal}$ while
$\mathsf{P}_{yy} = \mathsf{P}_{zz} = \SI{-0.0409}{\giga\pascal}$, and the
off-diagonal components stay at zero. The slopes at zero strain are the
crystal's elastic stiffnesses, $C_{11} = -\partial\mathsf{P}_{xx}/\partial
\epsilon_{xx} = \SI{3.920}{\giga\pascal}$ and $C_{12} = -\partial
\mathsf{P}_{yy}/\partial\epsilon_{xx} = \SI{2.324}{\giga\pascal}$. An even
squeeze applies all three strains at once; by the cubic symmetry
$\mathsf{P}_{xx}$ responds to $\epsilon_{yy}$ and $\epsilon_{zz}$ as
$\mathsf{P}_{yy}$ does to $\epsilon_{xx}$, so each diagonal component
changes by $-(C_{11} + 2C_{12})e$ while $\dd V/V = 3e$, and the bulk
modulus is $(C_{11} + 2C_{12})/3 = \SI{2.856}{\giga\pascal}$, the value
\cref{sec:pr-periodic} found from the slope of $P$.

ASE, like many libraries, reports the stress, which is $-\mathsf{P}$,
positive when the material is pulled; a program's sign convention is
worth checking before its numbers are used. ASE's \texttt{get\_stress}
returns it in \si{\electronvolt\per\angstrom\cubed} as six numbers,
\begin{equation*}
   -\mathsf{P}_{xx},\ -\mathsf{P}_{yy},\ -\mathsf{P}_{zz},\
   -\mathsf{P}_{yz},\ -\mathsf{P}_{xz},\ -\mathsf{P}_{xy} ,
\end{equation*}
and leaves out the kinetic part unless
it is called with \texttt{include\_ideal\_gas=True}; the pressure is then
minus the mean of the first three.

\begin{bridge}
In a density-functional calculation the stress is the same derivative of
the energy with respect to strain, and Nielsen and Martin showed how to
compute it from the electronic structure directly\cite{Nielsen1985}. A
basis of plane waves cut off at a fixed energy changes when the cell
does, so a cell that grows during a calculation is described by a basis
of different quality from the one it started with, and the stress
carries a spurious part, the Pulay stress; it is removed by raising the
cutoff until the stress no longer changes, or by correcting for the
incomplete basis as Francis and Payne set out\cite{Francis1990}. A basis
of orbitals centred on the atoms moves with them instead, and its
derivatives enter the stress as Pulay terms, as they enter the forces
(\cref{sec:en-atoms}).
\end{bridge}

\begin{keyidea}
The pressure tensor gives the force per unit area on a face of any
orientation, $\vb{F} = A\,\mathsf{P}\vb{n}$; it is $(\sum m v_\alpha v_\beta
+ \mathcal{W}_{\alpha\beta})/V$, from the response of the free energy to a
strain, and the mean of its diagonal is the pressure. A liquid's is $P$
times the unit matrix; a strained crystal's is not. Libraries report the
stress, $-\mathsf{P}$.
\end{keyidea}

\notebook{14}{stretch the crystal along one axis and read the tensor}

\begin{selfcheck}
ASE's stress, with the kinetic part included, is
\SI{0.002}{\electronvolt\per\angstrom\cubed} for each of its first three
numbers and zero for the others. What is the pressure in GPa, and is the
material pushing on its cell or pulling on it?
\end{selfcheck}
